\documentclass[review,nonacm,acmsmall]{acmart}

\setcopyright{none}
\citestyle{acmauthoryear}

\usepackage{booktabs}
\usepackage{graphicx}
\usepackage{pgfgantt}
\usepackage{float}
\usepackage{cleveref}
\usepackage{listings}
\usepackage{xcolor}
\usepackage[T1]{fontenc}
\usepackage[utf8]{inputenc}

\crefname{lstlisting}{listing}{listings}
\Crefname{lstlisting}{Listing}{Listings}

% Template-compatible optional Lean language definition.
% This paper uses Java, but the third template file is retained.
\definecolor{keywordcolor}{rgb}{0.7,0.1,0.1}
\definecolor{tacticcolor}{rgb}{0.0,0.1,0.6}
\definecolor{commentcolor}{rgb}{0.4,0.4,0.4}
\definecolor{symbolcolor}{rgb}{0.0,0.1,0.6}
\definecolor{sortcolor}{rgb}{0.1,0.5,0.1}
\definecolor{attributecolor}{rgb}{0.7,0.1,0.1}
% Optional Lean language listing definition retained from the course template.
% The current paper's lstset uses Java, not Lean.
% To list Lean code, use \begin{lstlisting}[language=lean] ... \end{lstlisting}.
\lstdefinelanguage{lean}{
  sensitive=true,
  morekeywords=[1]{import,prelude,protected,private,noncomputable,
    definition,lemma,variable,variables,theorem,example,open,as,
    export,axiom,inductive,structure,universe,universes,match,
    infix,infixl,infixr,notation,instance,eval,check,namespace,
    section,attribute,local,set_option,class,calc,have,show,by,
    in,at,let,forall,fun,exists,if,then,else,obtain,from,
    mutual,do,def,partial,mut,where,macro,syntax,deriving,
    return,try,catch,for,macro_rules,abbrev},
  morekeywords=[2]{Sort,Type,Prop},
  morekeywords=[3]{assumption,apply,intro,intros,allGoals,
    generalize,clear,revert,done,exact,refine,repeat,cases,
    rewrite,rw,simp,simp_all,contradiction,constructor,
    injection,induction},
  literate={:=}{{\color{symbolcolor}:=}}2
           {→}{{\ensuremath{\rightarrow}}}1
           {←}{{\ensuremath{\leftarrow}}}1
           {∀}{{\ensuremath{\forall}}}1
           {∃}{{\ensuremath{\exists}}}1
           {≤}{{\ensuremath{\leq}}}1
           {≥}{{\ensuremath{\geq}}}1,
  morecomment=[s][\color{commentcolor}]{/-}{-/},
  morecomment=[l][\color{commentcolor}]{--},
  showstringspaces=false,
  keepspaces=true,
  tabsize=3,
  extendedchars=false,
  breaklines=true,
  breakatwhitespace=true,
  basicstyle=\ttfamily\small,
  captionpos=b,
  columns=[l]fullflexible,
  identifierstyle={\ttfamily\color{black}},
  keywordstyle=[1]{\ttfamily\color{keywordcolor}},
  keywordstyle=[2]{\ttfamily\color{sortcolor}},
  keywordstyle=[3]{\ttfamily\color{tacticcolor}},
  stringstyle=\ttfamily,
  commentstyle={\ttfamily\footnotesize}
}


\title{Analysis-Guided Metamorphic Testing of Java Static Analyzers}

\author{Abdullah Al Abid (s260128)}
\author{Afonso Fonseca Morais de Sousa Pina (s263162)}
\author{Hasan Ahmed (s262722)}
\author{Mahmudul Hasan Aunoy (s260125)}
\author{Shohan Morol (s262724)}

% The project analyzes Java, so Java is selected instead of Lean.
\lstset{
    language=Java,
    basicstyle=\ttfamily\footnotesize,
    keywordstyle=\color{blue!65!black},
    columns=fullflexible,
    keepspaces=true,
    showstringspaces=false,
    breaklines=true,
    frame=single
}

\begin{document}
\maketitle

\section{Introduction}

\paragraph{Context}
Static program analysis predicts potential runtime errors without executing
every possible input. Yet the precision of an analyzer can depend on how a
program is expressed. Consider the following equivalent Java methods, where
the second reverses the branch condition and order:

\begin{minipage}[t]{0.47\linewidth}
\begin{lstlisting}
int f(int x) {
    if (x > 0)
        return 100 / x;
    return 0;
}
\end{lstlisting}
\end{minipage}
\hfill
\begin{minipage}[t]{0.47\linewidth}
\begin{lstlisting}
int f(int x) {
    if (x <= 0)
        return 0;
    return 100 / x;
}
\end{lstlisting}
\end{minipage}

For every Java integer input, both methods return the same value without
dividing by zero. However, an imprecise analyzer might warn on one
representation but not the other. Such a discrepancy does not necessarily
prove unsoundness: it can also reflect different levels of conservative
precision, which developers experience as warnings that appear or disappear
after harmless refactorings.

\paragraph{Gap}
Earlier DTU work applied semantics-preserving variations, such as renaming
variables and methods, to compare the robustness of an LLM and the Infer
analyzer in detecting null pointer exceptions, measuring results by reported
line numbers~\citep{chand2024}. The variations served as a robustness check
and were not placed using information about the program. Other work used
abstract interpretation to guide a fuzzer towards error-revealing inputs on
JPAMB~\citep{domingues2025}. However, a static analyzer's prediction for a
method does not depend on runtime inputs, so input-level guidance cannot
reveal whether the analyzer's answer depends on how the code is written.
That requires varying the program itself,
and it remains open where to vary it: rewriting at arbitrary sites may spend
effort on code whose analysis result cannot change.

\paragraph{Innovation}
We ask: \emph{can the abstract state computed by a lightweight analysis select
semantics-preserving rewrite sites that reveal prediction inconsistencies in a
Java static analyzer more efficiently than random or purely syntactic
selection?} Our key observation is that a rewrite can only change an
analyzer's result where the analyzer's precision depends on the \emph{form} of
the code: how a branch condition is written, where control flow merges, or how
a loop is laid out. For example, our sign analyzer does not narrow the
compared value at a branch, so it reports a possible division by zero in
\emph{both} versions of \texttt{f} above, and inverting that branch reveals
nothing. We expect that an analyzer which does narrow at branches can lose
that information when a rewrite moves where branches merge, so that the two
versions get different answers; milestone M1 tests this before we commit to
the full set of rewrites.

We therefore call a site \emph{uncertain} when, in the abstract state at that
program point, (i) a branch condition can still be both true and false, or
(ii) a divisor or array index may still be zero or out of bounds. A site is a
\emph{candidate} when it is uncertain \emph{and} one of our rewrites changes a
feature the analyzer's precision depends on there: branch polarity, where
branches merge, whether a jump goes backward, or the form of a comparison.
Rewrites are limited by a budget of $N$ rewrites for the whole benchmark, not
per method, because most JPAMB methods contain only one or two eligible sites.
The guided strategy spends this budget on candidate sites across all methods
first, choosing among them at random, and uses other eligible sites only when
no candidates remain; we expect rewrites at non-candidate sites to leave the
analyzer's answer unchanged and so waste budget. An \emph{inconsistency} is a
difference in the set of outcomes (\texttt{ok}, divide by zero, assertion
error, out of bounds, null pointer, \texttt{*}) predicted by the \emph{same}
analyzer for an original method and its rewritten version.

We will implement a parser-based site finder; a set of rewrites (branch
inversion with branch swapping, equivalent integer comparisons, loop
reshaping); the guided selector; and an interval domain alongside the existing
sign domain. Each rewrite has a precondition under which it preserves
behavior, which the site finder checks before applying it: for example, the
condition has no side effects, compares \texttt{int} values, and no rule
changes arithmetic that may overflow. To catch errors in the rules or their
implementation, we synthesize test inputs from the abstract state: for each
rewritten condition, inputs that make it true and false, plus the boundary
values of each variable's interval. We run the original and the variant on
these inputs and discard any variant whose outcome differs.

\paragraph{Evaluation}
We test three JPAMB analyzers that differ only in precision: sign analysis
without narrowing at branches, sign analysis with narrowing, and interval
analysis. They differ for reasons unrelated to our selection strategies, so
differences between strategies come from where they rewrite, not from how the
analyzers were built. Inconsistencies are always measured \emph{within} one
analyzer (original vs.\ variant). We compare the guided strategy with two
baselines over the same eligible sites and rewrite rules: random selection,
which draws uniformly from all eligible sites, and syntax-based selection,
which prefers sites close to a division or assertion. The guided strategy
uses the abstract states of our interval analysis, the most precise of the
three, because a coarser guide would mark almost every site as uncertain.

We compare the strategies in two settings. With a fixed \emph{rewrite
budget}, each strategy produces the same number $N$ of variants across the
benchmark; this measures how well it chooses sites. With a fixed \emph{time
budget}, each strategy runs for the same wall-clock time $T$ and produces as
many variants as it can; the time the guided strategy spends computing
abstract states counts towards $T$, so this measures whether the guidance is
worth its cost. Each setting is repeated with at least ten random seeds. We
measure distinct inconsistencies, valid variants (those that compile and pass
the synthesized tests), time to first inconsistency, and total runtime, and
compare strategies with a Mann--Whitney U test. Subjects are the integer and
branch methods of JPAMB plus a new set of cases that we submit as JPAMB
additions. We separately report rewrites that do not change the compiled
bytecode. Threats to validity are limited language coverage, correctness of
the rewrites, how representative the benchmark is, and that the guiding
analysis is also one of the analyzers under test.

\section{Plan}
\textbf{M1 (7 Nov):} a proof of concept on five method pairs with two
rewrites and a first version of branch narrowing in the sign analyzer, run
through \texttt{jpamb build} and \texttt{jpamb analyse}.
\textbf{M2 (14 Nov):} the site finder, the rewrites, the three selection
strategies, the interval domain, and full narrowing at branches.
\textbf{M3 (22 Nov):} an extended benchmark, and rewrites checked with test
inputs synthesized from the abstract state.
\textbf{M4 (29 Nov):} experiments, statistics, and the paper.
If M1 shows no inconsistencies for a rewrite, we replace that rewrite before
M2. If the interval domain is not ready at M2, development of the guided
strategy continues with the sign analysis with narrowing as a temporary
guide. Ownership is given in Table~\ref{tab:contributions} and the schedule
in Figure~\ref{fig:timeline}.

\clearpage
% % References section
\bibliographystyle{ACM-Reference-Format}
\bibliography{bibfile}

\section*{Project Timeline and Contributions}

\begin{figure}[H]
\centering
\caption{Proposed Project Timeline (2026)}
\label{fig:timeline}
\begin{ganttchart}[
    hgrid,
    time slot format=isodate,
    today=2026-10-21,
    today rule/.style={draw=blue!70!black, thick},
    today label font=\scriptsize,
    x unit=0.11cm,
    bar label node/.append style={text width=4.35cm, align=right},
    vgrid={*5{white},*2{dotted}},
    y unit chart=0.55cm,
    bar label font=\scriptsize,
    title label font=\small\bfseries,
    vrule/.style={very thick, red},
    vrule label font=\tiny\bfseries
  ]{2026-10-19}{2026-12-15}
  \gantttitlecalendar{month=shortname} \\
  \ganttbar{M1: Proof of concept and branch narrowing}{2026-10-28}{2026-11-07} \\
  \ganttbar{Semantics-preserving rewrites}{2026-10-30}{2026-11-14} \\
  \ganttbar{Interval domain and branch narrowing}{2026-11-01}{2026-11-14} \\
  \ganttbar{Random, syntactic and guided selection}{2026-11-01}{2026-11-14} \\
  \ganttbar{M3: Benchmark and test synthesis}{2026-11-08}{2026-11-22} \\
  \ganttbar{Experiments and statistical evaluation}{2026-11-18}{2026-11-27} \\
  \ganttbar{Related work and paper writing}{2026-11-09}{2026-11-29}
  \ganttvrule[vrule label node/.append style={anchor=north west}]{Proposal*}{2026-10-22}
  \ganttvrule{Paper*}{2026-11-29}
  \ganttvrule{Present*}{2026-12-06}
\end{ganttchart}

\vspace{2pt}

\end{figure}

\begin{table}[H]
\centering
\caption{Proposed technical contributions. $S_1$: Abdullah Al Abid (parsing
and rewriting); $S_2$: Hasan Ahmed (abstract domains and guidance);
$S_3$: Afonso Fonseca Morais de Sousa Pina (related work and test synthesis
for rewrite checking); $S_4$: Mahmudul Hasan Aunoy (JPAMB integration,
benchmark additions, syntactic analysis of the benchmark methods, and the
syntax-based baseline); $S_5$: Shohan Morol (experiments, statistics, and
part of the related work).}
\label{tab:contributions}
\scriptsize
\setlength{\tabcolsep}{3pt}
\begin{tabular}{@{}llrrrrrr@{}}
\toprule
ID & Contribution & Max & $S_1$ & $S_2$ & $S_3$ & $S_4$ & $S_5$ \\
\midrule
ISY & Syntactic analysis & 10 & 6 & 0 & 0 & 4 & 0 \\
NCR & Analysis-informed rewriting & 10 & 4 & 6 & 0 & 0 & 0 \\
NTS & Analysis-informed test synthesis & 10 & 0 & 5 & 5 & 0 & 0 \\
EXB & Accepted JPAMB additions & 5 & 0 & 0 & 0 & 5 & 0 \\
EST & Statistical evaluation & 5 & 0 & 0 & 0 & 0 & 5 \\
PRW & Relate recent papers & 10 & 0 & 0 & 7 & 0 & 3 \\
PWP & Write final paper & 10 & 2 & 2 & 2 & 2 & 2 \\
PEX & Illustrating examples & 4 & 0 & 0 & 2 & 2 & 0 \\
PDI & Process diagram & 3 & 0 & 0 & 0 & 3 & 0 \\
PMG & Project management & 6 & 3 & 3 & 0 & 0 & 0 \\
\midrule
& \textbf{Total} & \textbf{73} & 15 & 16 & 16 & 16 & 10 \\
& \textbf{Suggested maximum grade} & -- & 10 & 10 & 10 & 10 & 7 \\

\bottomrule
\end{tabular}
\end{table}

\end{document}
